\section{Experiments}

\subsection{Baselines}
We select six Stable Diffusion based methods as the baselines to our work. We categorize these baselines into three categories based on the cost to use them: (1) methods that add \textbf{no extra cost} to the pretrained model which includes we choose two versions of Stable Diffusion (SD), 1.4 and 2.1, as well as StrutureDiffusion~\citep{structurediffusion}, (2) methods that \textbf{increase inference time} which includes Composable Diffusion~\citep{composablediffusion} and MultiDiffusion (MD)~\citep{multidiffusion}, and (3) methods that \textbf{require additional training} for which we choose GLIGEN~\citep{gligen} as the baseline. Methodologies of the baselines are discussed in Section~\ref{related_works}. In the next sections, we group the results by the tree categories we introduce in this section. We include the implementation details in the appendix.





\subsection{Localized Text-to-Image Generation}
As mentioned in the introduction, previous works that tackle similar tasks, especially the ones with open vocabulary setting, usually rely on qualitative results and human judgements for evaluations. Due to human involvement in these evaluations, they tend to be expensive and inefficient, and are usually not scalable and prone to high variance. On the other hand, foundation models developed for many general purpose recognition tasks have shown great potential in supporting human to streamline various language and vision tasks. 
Here we investigate ways to integrate several different off-the-shelf large pretrained models for scalable automatic metrics for localized and compositional text-to-image generation. In each section below, we show that these large pretrained models have the ability to reflect correctly on the relative performance among the generations with different types of input information and agree with qualitative results and human evaluations.
\subsubsection{Generating with Bounding Boxes}
\begin{figure}[t]
  \centering
  \includegraphics[width=\textwidth]{img/coco.pdf}
  \caption{Illustration of generated images based on COCO bounding boxes.}
  \label{fig:coco}
\end{figure}
\paragraph{Experiment Setting and Dataset}
We use the validation set of COCO2017~\citep{coco} to perform bounding box based generation experiments. Each data point contains a caption for the overall scene and a set of bounding boxes each associated with a class label. Following the settings of~\citet{multidiffusion}, we create the pseudo text prompt with the class name for each bounding box and filter out 1095 examples with 2 to 4 non-human objects that are larger than 5\% of the image area. For models that are unable to take location information as inputs, we create pseudo text prompts by the format ``<caption> with <object1>, <object2>, ...''. In addition to all the baselines mentioned above, we also test the ability of MultiDiffusion as a plugin to models other than Stable Diffusion to compare with our method.
\paragraph{Evaluation Metrics}
We evaluate the generated images by (1) how close the generated image resembles a real COCO image (fidelity), (2) how consistent the generations are with the bounding boxes (controllability) and (3) how long it takes for the model to generate an image (inference time). We use Kernel Inception Score (KID)~\citep{kid} to evaluate fidelity. 
For controllability, we use an YOLOv8 model~\cite{yolo} trained on the COCO dataset to predict the bounding boxes in the generated images and then calculate the precision (P), recall (R), mAP50
and mAP50-95 
with the ground truth boxes. 
We use the default thresholds for all metrics. We also report the average inference time with a 50-step sampler on one NVIDIA Tesla V100 machine.
\paragraph{Results} 
Table~\ref{tab:bbox_results} shows the quantitative results of generation with bounding boxes. As we can observe, CAC improves the localized generation ability for all models we apply it to. As a plugin, CAC also does not substantially increase the inference time compared to MultiDiffusion. In Figure~\ref{fig:coco}, we illustrate a few qualitative examples that demonstrate the effect of CAC in Stable Diffusion (SD), MultiDiffusion (MD) and GLIGEN. For models that do not have localization ability like SD, CAC provides the new ability to generate contents based on the location information. For models that have localization ability such as MD and GLIGEN, CAC improves the generation by making the generated objects and features more recognizable. Notice that the performance is still strongly influenced by the base model. For example, both SD and SD+CAC generation are missing the ``woman'' in the caption.
\begin{figure}[t!]
  \centering
  \includegraphics[width=\textwidth]{img/cityscapes.pdf}
  \caption{Illustration of different approaches generating images via Cityscapes segmentation maps.}
  \label{fig:cityscapes}
\end{figure}

\begin{table}[t]
\caption{\small Experiment results with bounding box information.}
\label{tab:bbox_results}
\centering
\small
\begin{tabular}{@{}ccccccc@{}}
\toprule
\textbf{Method}                     & \textbf{KID $\downarrow$}     & \textbf{P $\uparrow$} & \textbf{R $\uparrow$} & \textbf{mAP50 $\uparrow$}  & \textbf{mAP50-95 $\uparrow$} & \textbf{Time (s) $\downarrow$}\\ \midrule
Ground Truth                        & -                & 0.6010             & 0.5680          & 0.6380          & 0.5460    & -            \\ \midrule
Stable Diffusion (SD) 1.4                & 0.00697          & 0.1120             & 0.0968          & 0.0452          & 0.0146      & 10.19      \\
Stable Diffusion (SD) 2.1                & 0.00733          & 0.0944             & 0.1050          & 0.0588          & 0.0178      & 9.20      \\
StructureDiffusion                  & \textbf{0.00654} & 0.0810             & 0.1070          & 0.0462          & 0.0147      &   \textbf{8.91}      \\ 
\textbf{SD 2.1 + CAC (Ours)} & 0.00786          & \textbf{0.2570}    & \textbf{0.1990} & \textbf{0.1650} & \textbf{0.0500}   &   9.28   \\ \midrule
Composable Diffusion                & 0.01174          & 0.1560             & 0.0852          & 0.0534          & 0.0165       & 32.75     \\
MultiDiffusion (MD)                      & 0.01189          & \textbf{0.3790}    & 0.2820          & 0.2570          & 0.1090      & 27.60      \\
\textbf{MD + CAC (Ours)}       & \textbf{0.00988} & \textbf{0.3790}    & \textbf{0.3050} & \textbf{0.2930} & \textbf{0.1340}  & \textbf{16.44} \\ \midrule
GLIGEN                              & 0.00691          & 0.7380             & 0.6280          & 0.6740          & 0.4670       & \textbf{27.60}     \\
GLIGEN + MD                         & \textbf{0.00679} & 0.6940             & 0.6530          & 0.6800          & 0.4690       & 101.15     \\
\textbf{GLIGEN + CAC (Ours)}               & 0.00708          & \textbf{0.7970}    & \textbf{0.7000} & \textbf{0.7810} & \textbf{0.5760}  & 27.85 \\ \bottomrule
\end{tabular}
\end{table}

\subsubsection{Generating with Semantic Segmentation Maps}
\paragraph{Experiment Setting and Dataset}
We use the validation set of Cityscapes~\citep{cityscapes} dataset for semantic segmentation map based generation. The dataset consists of 500 street photos taken in three cities and the pixel level semantic labeling from 30 predefined classes. We generate a pseudo caption for each image with the format ``a street photo in <city>'' where <city> indicates the city where the picture was taken. We also produce pseudo prompts for each semantic segmentation mask associated with each class by using the class name, and we use the same format to create prompts for models without localization ability. We also filtered all classes that occupy less than 5\% of the image. We center crop and resize each map to $512\times 512$ in order to match the dimensionality of the models. We omit the experiments related to GLIGEN due to the unavailability of their pretrained models on Cityscapes.
\paragraph{Evaluation Metric}
We also evaluate the model performance based on fidelity, controllability and inference time in this experiment. We use KID again to evaluate the fidelity. Since the divergence between data distribution of Cityscapes images and the generated images is substantial according to the KID results in Table~\ref{tab:semantic_seg_results}, semantic segmentation models that are only trained on Cityscapes will not work well in our setting due to the data distribution shift. As a result, we use Semantic Segment Anything (SSA)~\citep{ssa}, which is a general purpose open-vocabulary model that leverages Segment Anything Model (SAM)~\citep{sam} for semantic segmentation. We report the mean IoU (mIoU) score and mean accuracy (mACC) calculated from all classes, and all pixel accuracy (aACC).
\paragraph{Results}
Table~\ref{tab:semantic_seg_results} and Figure~\ref{fig:cityscapes} demonstrate the quantitative and qualitative results for the task of generation with semantic segmentation maps. Although the SSA model can achieve very high accuracy and IoU score for ground truth images, there is still a gap between performance on generated images and real images. Nevertheless, we still find the quantitative results resemble the relative performance of qualitative examination. In particular, like the previous experiments with bounding box information, our method can also provide additional localization ability with semantic segmentation maps. For localized method MD, the generated parts are usually well separated by not coherent, and CAC is able to create more consistent images compared to MD.
\begin{figure}[t]
\CenterFloatBoxes
\begin{floatrow}
\capttabbox{
\centering

\resizebox{0.45\textwidth}{!}{%
\begin{tabular}{@{}ccccc@{}}
\toprule
\textbf{Method}                            & \textbf{KID $\downarrow$}         & \textbf{aACC $\uparrow$}        & \textbf{mIoU $\uparrow$}        & \textbf{mACC $\uparrow$}         \\ \midrule
Ground Truth                               & - & 91.60                & 60.79                & 71.63                                                         \\ \midrule
Stable Diffusion 1.4                       & 0.128              & 34.85                & 4.54                 & 8.83                                                             \\
Stable Diffusion 2.1                       & \textbf{0.113}     & 36.37                & 4.55                 & 8.58                                                             \\
StructureDiffusion                             & 0.132              & 34.99                & 4.60                 & 8.82                                                           \\
\textbf{SD 2.1 + CAC (Ours)} & 0.115              & \textbf{51.29}       & \textbf{8.20}        & \textbf{13.20}                                                 \\ \midrule
Composable Diffusion                       & 0.166 & 47.03 & 6.13 & 10.05                                             \\
MultiDiffusion                             & 0.151              & \textbf{47.52}       & 8.60                 & 13.39                                                         \\
\textbf{MD + CAC (Ours)}       & \textbf{0.145}     & 46.13                & \textbf{8.61}        & \textbf{13.66}                                                 \\ \bottomrule
\end{tabular}
}
}{%
  \caption{\small Experiment results with semantic segmentation information.}
\label{tab:semantic_seg_results}
}

\ffigbox[\FBwidth]{%
  \includegraphics[width=0.45\textwidth]{img/ablation.pdf}
}{%
  
  \caption{\small Ablation study on the fideility-controllability tradeoff with and without CAC.}
  \label{fig:ablation}
}
\end{floatrow}
\end{figure}
\subsubsection{Generating with Compositional Prompts}
\begin{figure}[t!]
  \centering
  \includegraphics[width=\textwidth]{img/cc5002.pdf}
  \caption{\small Different generations with caption ``a blue backpack and a red chair'' from CC-500 dataset.}
\end{figure}
The challenge of localized generation stems not only from the incorporation of location information but also from the inherent difficulty of compositional generation. Studies in~\citet{winoground} show that compositionality, which is the task of creating complex content from combinations of simpler features and objects, is still an extremely challenging concept even in the era of large pretrained models.~\citet{composablediffusion, structurediffusion} also find that Stable Diffusion does not exempt from encountering this difficulty - different objects and the relationships between objects and attributes are frequently misrepresented or entirely missing when generating complex scenes.
Hence, in this section, we also investigate how our method can improve the performance of compositional generation.
\paragraph{Experiment Setting and Dataset}
We use the CC-500 dataset~\citep{structurediffusion} for this task. CC-500 consists of text prompts with the format ``a <color 1> <object 1> and a <color 2> <object 2>''. For methods that require location inputs, we apply a simple heuristic where the method will generate the first object on the left hand side of the image and the second object on the right hand side.
\paragraph{Evaluation Metric} 
In this task we mainly evaluate the accuracy of the generation. Following~\citet{structurediffusion}, we categorize each generation into three categories: generations that have incorrect or missing objects, generations that have the correct objects but the wrong colors, and the generations that have the correct objects with the correct colors. Similar to~\citet{gligen,structurediffusion}, we choose GLIP~\citep{glip}, the semantic rich open-vocabulary object detector to produce object and color detection results. We report the mean percentage of GLIP detection (mGLIP) with confidence threshold from 0.6 to 0.8 with 0.05 step size. Considering the difficulty in recognizing compositional objects, we also perform human evaluation on Amazon Mechanical Turk (mTurk) to verify the automatic results. Details about the human evaluation is in the appendix.
\paragraph{Results}
Table~\ref{tab:cc500_results} shows the quantitative results by automatic metrics and human evaluation. Even though there are some discrepancies between machine recognizability and human recognizability, both automatic evaluation and human annotators agree on the relative performance of the models and concur that our method significantly improves the compositional generation capability for all models. By localizing the features and the objects, our method is capable of creating better association between the attributes and the objects and render the generation more recognizable.
We also note that with the updated text encoder, SD 2.1 is able to handle compositional prompts remarkbly better than 1.4. 
Notice that while GLIGEN performs extremely well in the COCO setting, its performance in this experiment drastically drops. We hypothesize that this is because GLIGEN requires more complicated heuristics that represent the geometry of the scene correctly in order to generate more human recognizable objects. This performance gap reflects a fidelity-controllability tradeoff caused by the implementation of self attention control, which we discuss in the next ablation study section.
\begin{table}[t]
\caption{\small Experiment results on CC-500 dataset. The row with the highest "Correct Objects \& Correct Colors" rate is highlighted in each category.}
\label{tab:cc500_results}
\centering
\small
\begin{tabular}{@{}ccccccc@{}}
\toprule
\multirow{2}{*}{\textbf{Method}} & \multicolumn{2}{c}{\textbf{\begin{tabular}[c]{@{}c@{}}Missing/Incorrect\\ Object(s) $\downarrow$\end{tabular}}} & \multicolumn{2}{c}{\textbf{\begin{tabular}[c]{@{}c@{}}Correct Objects but\\ Wrong Color(s) $\downarrow$\end{tabular}}} & \multicolumn{2}{c}{\textbf{\begin{tabular}[c]{@{}c@{}}Correct Objects \&\\ Correct Colors $\uparrow$\end{tabular}}} \\ \cmidrule(l){2-7} 
                                 & \multicolumn{1}{c}{mGLIP}                  & \multicolumn{1}{c}{MTurk}                 & \multicolumn{1}{c}{mGLIP}                     & \multicolumn{1}{c}{MTurk}                     & \multicolumn{1}{c}{mGLIP}                     & \multicolumn{1}{c}{MTurk}                    \\ \midrule
Stable Diffusion (SD) 1.4        & 55.37\%                                          & 23.38\%                                         & 13.89\%                                             & 46.99\%                                             & 30.74\%                                             & 29.63\%                                            \\
Stable Diffusion (SD) 2.1        & 33.56\%                                          & 28.24\%                                         & 16.53\%                                             & 26.62\%                                             & 49.91\%                                             & 44.91\%                                            \\
StructureDiffusion               & 52.64\%                                          & 25.46\%                                         & 11.39\%                                             & 40.28\%                                             & 35.97\%                                             & 34.26\%                                            \\ 
\textbf{SD 2.1 + CAC (Ours)}     & \textbf{20.09\%}                                 & \textbf{7.41\%}                                 & \textbf{12.64\%}                                    & \textbf{11.34\%}                                    & \textbf{67.27\%}                                    & \textbf{81.25\%}                                   \\ \midrule
Composable Diffusion             & 51.25\%                                          & 21.99\%                                         & 13.94\%                                             & 34.03\%                                             & 34.81\%                                             & 43.75\%                                            \\
MultiDiffusion (MD)              & 20.05\%                                          & 10.88\%                                         & 16.48\%                                             & 9.72\%                                              & 63.47\%                                             & 79.40\%                                            \\
\textbf{MD + CAC (Ours)}         & \textbf{18.33\%}                                 & \textbf{7.41\%}                                 & \textbf{12.08\%}                                    & \textbf{4.40\%}                                     & \textbf{69.58\%}                                    & \textbf{87.96\%}                                   \\ \midrule
GLIGEN                           & 47.73\%                                          & 43.29\%                                         & 11.34\%                                             & 29.17\%                                             & 40.93\%                                             & 27.55\%                                            \\
\textbf{GLIGEN + CAC (Ours)}     & \textbf{31.11\%}                                 & \textbf{8.80\%}                                 & \textbf{8.98\%}                                     & \textbf{15.05\%}                                    & \textbf{59.91\%}                                    & \textbf{76.16\%}                                   \\ \bottomrule
\end{tabular}
\end{table}
\subsection{Ablation Study}
Fidelity-controllability tradeoff has been witnessed in many controllable generation methods~\citep{sdedit,e4e}: when the algorithm can control the generation better, i.e. the generation is more consistent with the user input, the generation quality of the images will usually decrease. In this section, we discuss the effect of cross attention control and self attention control on this tradeoff.

We select MultiDiffusion (MD) as the approach for self attention control and perform the COCO bounding box based generation with different ratios of diffusion timesteps to apply MD. To demonstrate the effect of cross attention control, we compare two settings with and without CAC applied to the timesteps without MD. Without CAC, we use standard Stable Diffusion with additional text prompts from the localization input concatenated to the caption for fair comparison. Figure~\ref{fig:ablation} shows the performance with and without CAC at various MD ratios. With high self attention control, i.e. large MD ratios, the model can achieve high mAP50 scores which represent better consistency with the bounding boxes. However, as we control the self attention in more timesteps, the images become less realistic and thus the KID values drop. Meanwhile, with CAC applied, the model is able to achieve higher mAP50 scores with lower MD ratios and lower KID scores. This indicates that models with CAC has better fidelity-controllability tradeoff than the models without CAC in this task.

